% Insert from benchmark/AegisKV/paper; independent figures may be composed by LaTeX.

% M1 | motivation-branch-gap | trace_observation
\begin{figure}[t]
\centering
\includegraphics[width=\columnwidth]{../../../output/pdf/aegiskv-paper-story-20260908/motivation-branch-gap.pdf}
\caption{Empirical distribution of the maximum parallel-branch duration gap in each of 124 observed workflows with overlapping branches. The horizontal axis is logarithmic. These observed duration differences are not validated causal CPM slack.}
\label{fig:motivation-branch-gap}
\end{figure}

% M2 | motivation-exposed-recovery | analytical
\begin{figure}[t]
\centering
\includegraphics[width=\columnwidth]{../../../output/pdf/aegiskv-paper-story-20260908/motivation-exposed-recovery.pdf}
\caption{Analytical illustration of exposed recovery delay, max(recovery delay - slack, 0), for slack values of 0, 5 and 10 seconds. Only one branch is delayed in a fixed DAG; other durations and resource conditions remain fixed. This is not an observed failure experiment.}
\label{fig:motivation-exposed-recovery}
\end{figure}

% M3 | motivation-kv-capacity | analytical
\begin{figure}[t]
\centering
\includegraphics[width=\columnwidth]{../../../output/pdf/aegiskv-paper-story-20260908/motivation-kv-capacity.pdf}
\caption{Analytical KV footprint for Qwen3-8B BF16 at 147456 bytes per token (Ki=1024). At 96 Ki tokens, one and two copies occupy 13.5 and 27 GiB. The dotted line subtracts model weights and 2 GiB runtime from 40 GiB HBM, before dynamic foreground KV; it is not a lendable Peer quota.}
\label{fig:motivation-kv-capacity}
\end{figure}

% S1 | system-overview | design
\begin{figure*}[t]
\centering
\includegraphics[width=\textwidth]{../../../output/pdf/aegiskv-paper-story-20260908/system-overview.pdf}
\caption{AegisKV revision-3 design. Solid arrows carry KV data and dashed arrows carry control. The source writes independently to Peer HBM and Store DRAM; committed backup data must be materialized into a new primary for recovery. The complete design is broader than the implementation subset evaluated here.}
\label{fig:system-overview}
\end{figure*}

% S2 | design-slack-tickets | design
\begin{figure*}[t]
\centering
\includegraphics[width=\textwidth]{../../../output/pdf/aegiskv-paper-story-20260908/design-slack-tickets.pdf}
\caption{Algorithm 1 turns a residual DAG and remaining-time estimates into versioned slack tickets. ES and LS denote earliest and latest start times; m is a conservative margin. Unknown structure, missing estimates or stale versions invalidate affected tickets to zero. Predictor accuracy is not established by this design illustration.}
\label{fig:design-slack-tickets}
\end{figure*}

% S3 | design-shared-prefix | design
\begin{figure*}[t]
\centering
\includegraphics[width=\textwidth]{../../../output/pdf/aegiskv-paper-story-20260908/design-shared-prefix.pdf}
\caption{Shared-prefix accounting in Algorithm 2. Consumers A and B share physical chunks C0 and C1. Evicting C1 requires its committed descendants C2-A and C2-B to be evicted as a closure. Solid arrows are prefix-parent relationships, dashed arrows are consumer references, and shared physical bytes are counted once.}
\label{fig:design-shared-prefix}
\end{figure*}

% S4 | design-foreground-budget | design
\begin{figure*}[t]
\centering
\includegraphics[width=\textwidth]{../../../output/pdf/aegiskv-paper-story-20260908/design-foreground-budget.pdf}
\caption{Algorithm 3 checks memory and transfer envelopes before allocation, copying or quota growth. A fresh QoS certificate is additionally required. Growth needs stable headroom, demand and completed reclamation; invalid conditions revoke permission, and capacity remains charged until a free ACK. This illustrates the full design, not a certificate obtained by the simulation.}
\label{fig:design-foreground-budget}
\end{figure*}

% S5 | commit-tier-timeline | design
\begin{figure*}[t]
\centering
\includegraphics[width=\textwidth]{../../../output/pdf/aegiskv-paper-story-20260908/commit-tier-timeline.pdf}
\caption{Independent tier commits for one chunk. The Peer copy is acknowledged before source failure, while the pending DRAM copy aborts. Horizontal spacing denotes logical event order, not measured duration. A committed Peer copy does not imply a committed DRAM copy.}
\label{fig:commit-tier-timeline}
\end{figure*}

% S6 | commit-recoverable-prefix | design
\begin{figure*}[t]
\centering
\includegraphics[width=\textwidth]{../../../output/pdf/aegiskv-paper-story-20260908/commit-recoverable-prefix.pdf}
\caption{Recoverable prefix after the illustrated failure. Peer HBM has committed chunks 0-1, whereas Store DRAM has only chunk 0. Aborted or absent chunks provide no recovery tokens. Each tier computes its own continuous committed prefix; these are illustrative chunks.}
\label{fig:commit-recoverable-prefix}
\end{figure*}

% S7 | commit-materialize-recompute | design
\begin{figure*}[t]
\centering
\includegraphics[width=\textwidth]{../../../output/pdf/aegiskv-paper-story-20260908/commit-materialize-recompute.pdf}
\caption{A committed backup prefix must first be copied into new primary buffers. The unprotected suffix is then recomputed before decoding resumes. Arrows are conceptual and do not encode time or bytes; prefix materialization and decode resumption are distinct milestones.}
\label{fig:commit-materialize-recompute}
\end{figure*}

% E1 | workload-prompt_tokens | workload_description
\begin{figure}[t]
\centering
\includegraphics[width=\columnwidth]{../../../output/pdf/aegiskv-paper-story-20260908/workload-prompt_tokens.pdf}
\caption{Prompt-length distribution of the 1304 original requests in the frozen FanOutQA-derived workload (Ki=1024). Synthetic revisit requests are excluded. The maximum total context is 86214 tokens, so this workload does not establish coverage of 128K contexts.}
\label{fig:workload-prompt_tokens}
\end{figure}

% E2 | workload-workflow_call_counts | workload_description
\begin{figure}[t]
\centering
\includegraphics[width=\columnwidth]{../../../output/pdf/aegiskv-paper-story-20260908/workload-workflow_call_counts.pdf}
\caption{Distribution of original call counts across 160 frozen workflows, with a maximum of 27 calls. A sample is a workflow. These counts precede the controlled duplication of the complete DAG into two rounds.}
\label{fig:workload-workflow_call_counts}
\end{figure}

% E3 | f0-revisit-wct | modeled_exploration
\begin{figure*}[t]
\centering
\includegraphics[width=\textwidth]{../../../output/pdf/aegiskv-paper-story-20260908/f0-revisit-wct.pdf}
\caption{Mechanism simulation of the controlled two-round workload: 160 paired workflows, 2608 requests, C=16, 1 TB DRAM and a 25 GB/s shared path. One execution per system. WCT runs from first release to final request completion, excluding admission waiting outside active workflow slots. The same ordering by Mooncake's no-failure revisit WCT is used in every workflow view; each position denotes one workflow, and lines do not denote a time series. }
\label{fig:f0-revisit-wct}
\end{figure*}

% E4 | f0-revisit-wct-difference | modeled_exploration
\begin{figure*}[t]
\centering
\includegraphics[width=\textwidth]{../../../output/pdf/aegiskv-paper-story-20260908/f0-revisit-wct-difference.pdf}
\caption{Mechanism simulation of the controlled two-round workload: 160 paired workflows, 2608 requests, C=16, 1 TB DRAM and a 25 GB/s shared path. One execution per system. WCT difference from Mooncake for each paired workflow; positive values are slower. Mooncake is the zero reference. The same ordering by Mooncake's no-failure revisit WCT is used in every workflow view; each position denotes one workflow, and lines do not denote a time series. }
\label{fig:f0-revisit-wct-difference}
\end{figure*}

% E5 | f0-revisit-wct-cdf | modeled_exploration
\begin{figure}[t]
\centering
\includegraphics[width=\columnwidth]{../../../output/pdf/aegiskv-paper-story-20260908/f0-revisit-wct-cdf.pdf}
\caption{Mechanism simulation of the controlled two-round workload: 160 paired workflows, 2608 requests, C=16, 1 TB DRAM and a 25 GB/s shared path. One execution per system. Empirical WCT distributions contain all 160 workflows per system, including the full tail. Workflows describe within-execution heterogeneity, not independent execution replicates; no population confidence interval is inferred.}
\label{fig:f0-revisit-wct-cdf}
\end{figure}

% E6 | f0-revisit-dram-hit | modeled_exploration
\begin{figure*}[t]
\centering
\includegraphics[width=\textwidth]{../../../output/pdf/aegiskv-paper-story-20260908/f0-revisit-dram-hit.pdf}
\caption{Mechanism simulation of the controlled two-round workload: 160 paired workflows, 2608 requests, C=16, 1 TB DRAM and a 25 GB/s shared path. One execution per system. For each workflow, DRAM token hit rate is actually reused DRAM prompt tokens divided by all prompt tokens, including both cold and revisit rounds. Recovery tokens are excluded. The same ordering by Mooncake's no-failure revisit WCT is used in every workflow view; each position denotes one workflow, and lines do not denote a time series. }
\label{fig:f0-revisit-dram-hit}
\end{figure*}

% E7 | f0-revisit-cache-sources | modeled_exploration
\begin{figure*}[t]
\centering
\includegraphics[width=\textwidth]{../../../output/pdf/aegiskv-paper-story-20260908/f0-revisit-cache-sources.pdf}
\caption{Mechanism simulation of the controlled two-round workload: 160 paired workflows, 2608 requests, C=16, 1 TB DRAM and a 25 GB/s shared path. One execution per system. Exclusive decomposition of all prompt tokens into Local HBM reuse, Store DRAM reuse, Peer HBM reuse and computed tokens. Fractions sum to 100\%; normal Peer reuse is zero in these runs. Backup bytes are not counted as foreground cache hits.}
\label{fig:f0-revisit-cache-sources}
\end{figure*}

% E8 | f0-revisit-scheduler-queue | modeled_exploration
\begin{figure*}[t]
\centering
\includegraphics[width=\textwidth]{../../../output/pdf/aegiskv-paper-story-20260908/f0-revisit-scheduler-queue.pdf}
\caption{Mechanism simulation of the controlled two-round workload: 160 paired workflows, 2608 requests, C=16, 1 TB DRAM and a 25 GB/s shared path. One execution per system. Workflow means of request delay from actual GPU scheduler admission to first compute. CPU, GET and destination-HBM reservation delays are separate. The same ordering by Mooncake's no-failure revisit WCT is used in every workflow view; each position denotes one workflow, and lines do not denote a time series. }
\label{fig:f0-revisit-scheduler-queue}
\end{figure*}

% E9 | f0-revisit-hbm-wait | modeled_exploration
\begin{figure*}[t]
\centering
\includegraphics[width=\textwidth]{../../../output/pdf/aegiskv-paper-story-20260908/f0-revisit-hbm-wait.pdf}
\caption{Mechanism simulation of the controlled two-round workload: 160 paired workflows, 2608 requests, C=16, 1 TB DRAM and a 25 GB/s shared path. One execution per system. Workflow means of foreground DRAM GET destination-HBM reservation waiting, including zero-wait requests. GET populations differ with cache hits; cumulative wait is not an additive WCT penalty. The same ordering by Mooncake's no-failure revisit WCT is used in every workflow view; each position denotes one workflow, and lines do not denote a time series. }
\label{fig:f0-revisit-hbm-wait}
\end{figure*}

% E10 | f0-revisit-hbm-pressure | modeled_exploration
\begin{figure*}[t]
\centering
\includegraphics[width=\textwidth]{../../../output/pdf/aegiskv-paper-story-20260908/f0-revisit-hbm-pressure.pdf}
\caption{Mechanism simulation of the controlled two-round workload: 160 paired workflows, 2608 requests, C=16, 1 TB DRAM and a 25 GB/s shared path. One execution per system. For each of all eight GPUs, event-interval integration measures the fraction of workload time with less than 1 GiB unallocated HBM. Reclaimable caches remain allocated. Denominators are each system's workload makespan; this is not hardware HBM bandwidth utilization.}
\label{fig:f0-revisit-hbm-pressure}
\end{figure*}

% E11 | f0-revisit-ttft-cdf | modeled_exploration
\begin{figure}[t]
\centering
\includegraphics[width=\columnwidth]{../../../output/pdf/aegiskv-paper-story-20260908/f0-revisit-ttft-cdf.pdf}
\caption{Mechanism simulation of the controlled two-round workload: 160 paired workflows, 2608 requests, C=16, 1 TB DRAM and a 25 GB/s shared path. One execution per system. Empirical time-to-first-token distributions for all 2608 requests per system. The time axis is logarithmic and retains the complete range. These modeled values are not a native-serving QoS certificate.}
\label{fig:f0-revisit-ttft-cdf}
\end{figure}

% E12 | f0-revisit-tpot-tail | modeled_exploration
\begin{figure*}[t]
\centering
\includegraphics[width=\textwidth]{../../../output/pdf/aegiskv-paper-story-20260908/f0-revisit-tpot-tail.pdf}
\caption{Mechanism simulation of the controlled two-round workload: 160 paired workflows, 2608 requests, C=16, 1 TB DRAM and a 25 GB/s shared path. One execution per system. Inclusive empirical survival of request-mean time per output token (TPOT), for 2432 requests with defined TPOT per system. Both axes are logarithmic. Each observation is a request mean rather than an individual inter-token gap; extreme values are retained.}
\label{fig:f0-revisit-tpot-tail}
\end{figure*}

% E13 | f0-revisit-controller-service-difference | modeled_exploration
\begin{figure*}[t]
\centering
\includegraphics[width=\textwidth]{../../../output/pdf/aegiskv-paper-story-20260908/f0-revisit-controller-service-difference.pdf}
\caption{Mechanism simulation of the controlled two-round workload: 160 paired workflows, 2608 requests, C=16, 1 TB DRAM and a 25 GB/s shared path. One execution per system. Paired controller service demand (foreground lookup plus admission), relative to Mooncake. Service quotations are fixed; the 50 microsecond foreground lookup cost is assumed. This is modeled CPU service, not native CPU utilization. The same ordering by Mooncake's no-failure revisit WCT is used in every workflow view; each position denotes one workflow, and lines do not denote a time series. }
\label{fig:f0-revisit-controller-service-difference}
\end{figure*}

% E14 | f0-revisit-controller-breakdown | modeled_exploration
\begin{figure*}[t]
\centering
\includegraphics[width=\textwidth]{../../../output/pdf/aegiskv-paper-story-20260908/f0-revisit-controller-breakdown.pdf}
\caption{Mechanism simulation of the controlled two-round workload: 160 paired workflows, 2608 requests, C=16, 1 TB DRAM and a 25 GB/s shared path. One execution per system. Additive controller service demand separated into foreground lookup and admission. Totals are 2.58, 5.44 and 6.05 seconds for Mooncake, DRAM-only AegisKV and AegisKV. Waiting time is excluded; this sum is neither request latency nor CPU utilization.}
\label{fig:f0-revisit-controller-breakdown}
\end{figure*}

% E15 | f0-revisit-background-copy | modeled_exploration
\begin{figure*}[t]
\centering
\includegraphics[width=\textwidth]{../../../output/pdf/aegiskv-paper-story-20260908/f0-revisit-background-copy.pdf}
\caption{Mechanism simulation of the controlled two-round workload: 160 paired workflows, 2608 requests, C=16, 1 TB DRAM and a 25 GB/s shared path. One execution per system. Actual DRAM-write plus Peer-copy bytes, attributed to the source workflow, including partially completed aborted copies and cleanup traffic. Values are cumulative transfer volume, not resident capacity or instantaneous bandwidth. The same ordering by Mooncake's no-failure revisit WCT is used in every workflow view; each position denotes one workflow, and lines do not denote a time series. }
\label{fig:f0-revisit-background-copy}
\end{figure*}

% E16 | f0-revisit-peer-copy | modeled_exploration
\begin{figure*}[t]
\centering
\includegraphics[width=\textwidth]{../../../output/pdf/aegiskv-paper-story-20260908/f0-revisit-peer-copy.pdf}
\caption{Mechanism simulation of the controlled two-round workload: 160 paired workflows, 2608 requests, C=16, 1 TB DRAM and a 25 GB/s shared path. One execution per system. Peer-copy bytes are counted directly from Peer flows. Mooncake and DRAM-only AegisKV have overlapping zero lines because the path is disabled. Cross-system total-traffic subtraction is not used to estimate this component. The same ordering by Mooncake's no-failure revisit WCT is used in every workflow view; each position denotes one workflow, and lines do not denote a time series. }
\label{fig:f0-revisit-peer-copy}
\end{figure*}

% E17 | f0-revisit-traffic-breakdown | modeled_exploration
\begin{figure*}[t]
\centering
\includegraphics[width=\textwidth]{../../../output/pdf/aegiskv-paper-story-20260908/f0-revisit-traffic-breakdown.pdf}
\caption{Mechanism simulation of the controlled two-round workload: 160 paired workflows, 2608 requests, C=16, 1 TB DRAM and a 25 GB/s shared path. One execution per system. Transferred bytes are partitioned into foreground GET, DRAM writes and Peer copies, including transferred portions of aborted flows and cleanup. TiB denotes 2\textasciicircum{}40 bytes. DRAM-write volume also changes when Peer protection is enabled.}
\label{fig:f0-revisit-traffic-breakdown}
\end{figure*}

% E18 | f0-revisit-transfer-service_slowdown | modeled_exploration
\begin{figure}[t]
\centering
\includegraphics[width=\columnwidth]{../../../output/pdf/aegiskv-paper-story-20260908/f0-revisit-transfer-service_slowdown.pdf}
\caption{Mechanism simulation of the controlled two-round workload: 160 paired workflows, 2608 requests, C=16, 1 TB DRAM and a 25 GB/s shared path. One execution per system. Successful foreground GET data-service time divided by isolated bottleneck service time for the same bytes; fixed setup is excluded. There are 1284, 773 and 662 GETs respectively. Population and byte-size differences prevent an isolated causal interpretation of the CDF separation.}
\label{fig:f0-revisit-transfer-service_slowdown}
\end{figure}

% E19 | f1-q25-workflow-penalty | modeled_exploration
\begin{figure*}[t]
\centering
\includegraphics[width=\textwidth]{../../../output/pdf/aegiskv-paper-story-20260908/f1-q25-workflow-penalty.pdf}
\caption{Mechanism simulation on the original 160-workflow, 1304-request workload; C=16, 1 TB DRAM, 25 GB/s. Each system is compared with its own matched no-failure execution. A focal request is preselected nearest 25\% baseline workflow completion and failed at 60\% fresh prefill progress. All 160 signed workflow WCT differences are shown in the shared ordering used by the no-failure revisit figures; each system is executed once.}
\label{fig:f1-q25-workflow-penalty}
\end{figure*}

% E20 | f1-q50-workflow-penalty | modeled_exploration
\begin{figure*}[t]
\centering
\includegraphics[width=\textwidth]{../../../output/pdf/aegiskv-paper-story-20260908/f1-q50-workflow-penalty.pdf}
\caption{Mechanism simulation on the original 160-workflow, 1304-request workload; C=16, 1 TB DRAM, 25 GB/s. Each system is compared with its own matched no-failure execution. A focal request is preselected nearest 50\% baseline workflow completion and failed at 60\% fresh prefill progress. All 160 signed workflow WCT differences are shown in the shared ordering used by the no-failure revisit figures; each system is executed once.}
\label{fig:f1-q50-workflow-penalty}
\end{figure*}

% E21 | f1-q75-workflow-penalty | modeled_exploration
\begin{figure*}[t]
\centering
\includegraphics[width=\textwidth]{../../../output/pdf/aegiskv-paper-story-20260908/f1-q75-workflow-penalty.pdf}
\caption{Mechanism simulation on the original 160-workflow, 1304-request workload; C=16, 1 TB DRAM, 25 GB/s. Each system is compared with its own matched no-failure execution. A focal request is preselected nearest 75\% baseline workflow completion and failed at 60\% fresh prefill progress. All 160 signed workflow WCT differences are shown in the shared ordering used by the no-failure revisit figures; each system is executed once.}
\label{fig:f1-q75-workflow-penalty}
\end{figure*}

% E22 | failure-progress-p95-penalty | modeled_exploration
\begin{figure*}[t]
\centering
\includegraphics[width=\textwidth]{../../../output/pdf/aegiskv-paper-story-20260908/failure-progress-p95-penalty.pdf}
\caption{Mechanism simulation on the original 160-workflow, 1304-request workload; C=16, 1 TB DRAM, 25 GB/s. Each system is compared with its own matched no-failure execution. P95 of 160 paired workflow WCT differences at each predeclared fault location, including workflows completed before failure. The x axis is realized baseline completion fraction. Later locations include more already-completed workflows; this curve does not isolate recovery difficulty. Percentiles are taken after pairing, and the three interventions are not independent replicates.}
\label{fig:failure-progress-p95-penalty}
\end{figure*}

% E23 | failure-focal-recovery-phases | modeled_exploration
\begin{figure*}[t]
\centering
\includegraphics[width=\textwidth]{../../../output/pdf/aegiskv-paper-story-20260908/failure-focal-recovery-phases.pdf}
\caption{Mechanism simulation on the original 160-workflow, 1304-request workload; C=16, 1 TB DRAM, 25 GB/s. Each system is compared with its own matched no-failure execution. Recovery of the one preselected focal request at each location. The first phase runs from failure to continuous-prefix materialization, including reservation, setup and copying. The second runs to decode resumption, including suffix computation and scheduling. Phase boundaries come from distinct events; bars are not averages over all affected requests.}
\label{fig:failure-focal-recovery-phases}
\end{figure*}

% E24 | sensitivity-bandwidth-mean_wct_s | modeled_exploration
\begin{figure*}[t]
\centering
\includegraphics[width=\textwidth]{../../../output/pdf/aegiskv-paper-story-20260908/sensitivity-bandwidth-mean_wct_s.pdf}
\caption{Mean WCT in the two-round mechanism workload at 5, 25 and 50 GB/s; C=16 and DRAM=1 TB are fixed. Each point summarizes 160 workflows in one execution per system. The bandwidth axis is linear and segments only connect tested points; no independent-repeat confidence interval is inferred.}
\label{fig:sensitivity-bandwidth-mean_wct_s}
\end{figure*}

% E25 | sensitivity-concurrency-mean_wct_s | modeled_exploration
\begin{figure*}[t]
\centering
\includegraphics[width=\textwidth]{../../../output/pdf/aegiskv-paper-story-20260908/sensitivity-concurrency-mean_wct_s.pdf}
\caption{Mean WCT in the two-round mechanism workload at 8, 16 and 32 active workflows; DRAM=1 TB and bandwidth=25 GB/s are fixed. Each point is one execution with 160 workflows. This closed-loop concurrency sweep is not throughput at a fixed open-loop arrival rate.}
\label{fig:sensitivity-concurrency-mean_wct_s}
\end{figure*}

% E26 | sensitivity-dram-capacity-dram_hit_pct | modeled_exploration
\begin{figure*}[t]
\centering
\includegraphics[width=\textwidth]{../../../output/pdf/aegiskv-paper-story-20260908/sensitivity-dram-capacity-dram_hit_pct.pdf}
\caption{DRAM prompt-token hit rate in the two-round mechanism workload at 32, 128 and 931.3 GiB, using a base-2 logarithmic capacity axis. C=16 and bandwidth=25 GB/s are fixed. The denominator includes both cold and revisit prompt tokens; each point is one execution per system.}
\label{fig:sensitivity-dram-capacity-dram_hit_pct}
\end{figure*}
